\chapter{Introduction} \label{chapter introduction}

\section{Aim of this book}

This book presents an analysis of the Japanese DP as part of the research program \textit{Cartography} (\citealt{pollock:1989:verb, cinque:1993:evidence, rizzi:1997:the, cinque:1999:adverbs, cinque:2006:complement, cinque:2010:adj, cinque:2017:micro, cinque:2020:relcl, cinque:2023:linearization} among others). Its main goal is to derive a structure of the DP and all its domains in which all types of modifiers can be incorporated. This DP structure to be proposed is built on previous cartographic analyses of the DP \citep{cinque:1993:evidence, cinque:2010:adj, cinque:2020:relcl, svenonius:2008:position, laenzlinger:2005:french, laenzlinger:2011:elements, kim:2019:syntax}. Importantly, it does not only \textit{accommodate} the Japanese data, but, rather, is \textit{enriched} by it, shaping our understanding of the DP and nominal modification as a research topic.

Specifically, Cartography provides the necessary tool for an analysis of the Japanese DP, highlighting Japanese as a language which does not behave entirely differently from other languages analyzed as part of this research field. There are, however, significant points of deviation, which I argue to be even more important for future analyses of the DP. With the in-depth investigation conducted in this book, I hope to achieve two goals: To reveal new insights pertaining to the Japanese nominal domain and to allow for cross-linguistic generalizations from this language, which will enhance our linguistic knowledge and the inventory of our linguistic tools and categories.

In this introductory chapter, I will present the scope and the main research questions of this book, as well as the main motivation for writing it, followed by the most important claims. In \sectref{intro japanese}, I present an overview of the Japanese nominal modification system. In the final section, I provide an overview of how data was collected.


\section{Scope of this book and questions to be answered}

It is a characteristic of Japanese that the syntactic role a modifier performs within the DP is not visible from the surface structure, and neither are its syntactic function or its syntactic position. This can be attributed to two main morphosyntactic peculiarities. First, in most cases, there is no morphological difference between attributively used lexemes and their predicative counterparts. Second, and perhaps more importantly, there is \textit{only one} surface construction, only one mechanism available for lexemes partaking in nominal modification. Consider the minimal pairs below which present an \textit{i}-adjective \ref{doresuintro} and a verb \ref{karegahonwokatta} in predicative and in attributive position.\footnote{The verb is inflected for past tense, highlighted via the element -\textsc{ta}, or the allomorph /da/. See Section \ref{chapter 1, i and na} for a first discussion on the morpheme \textsc{i}.} Essentially, it seems as if the word order has been shifted around and we only know that the b) examples involve a modification structure from the very strict SOV order present in Japanese.

\ex. \label{doresuintro} \ag. Ano doresu-wa taka-i.\\
that dress-\textsc{top} expensive-\textsc{i}\\
`That dress is expensive.'
\bg. ano taka-i doresu\\
that expensive-\textsc{i} dress\\
`that expensive dress' \null \hfill \textit{i}-adjective, non-past

\ex. \label{karegahonwokatta} \ag. Kare-ga [hon-wo kat-ta].\\
he-\textsc{nom} book-\textsc{acc} buy-\textsc{pst}\\
`He bought a book.’
\bg. [kare-ga kat-ta] hon\\
he-\textsc{nom} buy-\textsc{pst} book\\
`the book (which) he bought' \null \hfill verb, past \label{kare-ga}

Now, compare \ref{doresuintro} to the different corresponding English DPs depicted below. \is{English}

\ex. \a. the expensive dress \label{expensive a}
\b. the dress which was expensive \label{expensive b}
\c. a dress, expensive but beautiful \label{expensive c}

The English adjective \textit{expensive}, besides performing regular adjectival modification in \ref{expensive a}, can appear inside a relative clause \ref{expensive b}. Such relative clauses are explicitly marked via relative pronouns or complementizers, elements which are absent in Japanese \citep{kuno:1973:structure, comrie:1998:rethinking}. Furthermore, the adjective can occur postnominally as depicted in \ref{expensive c}. Postnominal adjectives in English are typically analyzed as \is{reduced relative clause}\textit{reduced relative clauses} (RRCs; \citealt{sadler:1994:prenominal, larson:1998:events, cinque:2010:adj}). They are reduced since they lack a relative pronoun/complementizer and a copula/verb \citep{douglas:2016:syntactic, harwood:2018:rrc}. This reduction can be depicted as follows:

\ex. \a. a dress \st{which/that} \st{is} expensive but beautiful
\b. the stars \st{which/that} \st{are} visible

In many languages, adjectives agree with the noun in case, number and gender. In some languages, including \is{German} German and \is{Dutch} Dutch, they do this only in attributive \ref{rot a}, not in predicative position \ref{rot b}. This predicts agreement to be also absent when the adjective appears inside a relative clause \ref{rot c} (\citealt{roehrs:2008:something, belk:2017:attributes, zeijlstra:2020:agree} a.o.).

\ex. \ag. der rot-\textbf{e} Tisch (German)\\
the red-\textsc{nom.sg.m} table.\textsc{nom.sg.m}\\
`the red table’ \label{rot a}
\bg. Der Tisch ist rot.\\
the table.\textsc{nom.sg.m} is red.$\emptyset$\\
`The table is red.’ \label{rot b}
\cg. der Tisch, der rot ist\\
the table.\textsc{nom.sg.m} the red.$\emptyset$ is\\
`the table which is red’ \label{rot c}

Nominal modifiers with a more predicative structure are called \textit{indirect} or \textit{predicative} modifiers, contrasting them with \textit{direct} or \textit{attributive} modifiers, those not related to predicativity \citep{sproat:1991:the, alex:2007:nounphrase, cinque:2010:adj}. Another way in which languages can cue different types of modification is through the optional appearance of additional morphosyntactic elements. For example, the element \textsc{de} in \is{Mandarin} Mandarin Chinese, which is optional with adjectival modifiers, but mandatory with verbal modifiers, has been analyzed as a clear indicator of indirect modification \citep[571]{sproat:1991:the}.

\exg. h\v{a}o (de) tiānqì\\
good \textsc{de} weather\\
`good weather’ \null \hfill \citep[573]{sproat:1991:the}

In Japanese, however, where no agreement between head and modifier is observable, where modifiers can only appear prenominally, and where only one modification structure exists, a simple modifier in non-past as given above in \ref{doresuintro}, is structurally ambiguous. These facts have led to the prevailing assumption that this language simply lacks direct modification, with all modifiers being interpreted as relative clauses (for example \citealt{kuno:1973:structure, hinds:1988:japanese, whitman:1981:internal, kaplan:1995:the}).

Admittedly, structural ambiguity also holds for cases in \is{English} English in which a single adjectival modifier modifies a noun, such as \textit{red book}. However, by relating the two types of modifiers to two different structural positions, besides their two different internal structures, Cartography offers another diagnostic, namely ordering. Observe \ref{big red}. \is{adjective ordering restrictions}

\ex. \label{big red} \a. big red book
\b. ??red big book

One order of the adjectives in \ref{big red} is clearly more natural than the other, given an out-of-the-blue context. Rather than assuming cognitive principles, Cartography explains this behavior through selection. Simply said, the projection hosting direct modifiers of \textsc{size} structurally dominates the one hosting direct modifiers of \textsc{color}. Importantly, in the indirect domain, where several functional projections of the same category can co-exist, no structural dominance exists, predicting the absence of strict order between indirect modifiers of the same type \citep{sproat:1991:the, baker:2003:verbal}. \is{English}

\ex. \a. book that is big that is red
\b. book that is red that is big

Crucially, the claim that modification in Japanese involves only relative clauses extends to \is{absence of ordering restrictions}ordering restrictions, or rather the absence thereof. Different from the minimal pair in English given in \ref{big red}, no ordering restrictions, or even tendencies, exist in the corresponding Japanese example \ref{okiakaiakaioki}.

\ex. \label{okiakaiakaioki} \ag. ōki-i aka-i hon\\
big-\textsc{i} red-\textsc{i} book\\
\bg. aka-i ōki-i hon\\
red-\textsc{i} big-\textsc{i} book\\
`big red book’

Even if we were to maintain the null hypothesis that lexical items in every language, unless proven otherwise, can appear as both types of modifiers, direct or indirect, these facts named above make it hard to argue for Japanese as a language in which both types exist. This claim also persists in the few theoretical works that discuss this issue with regard to this language \citep{sproat:1991:the, baker:2003:verbal, laenzlinger:2011:elements}.

Apart from the fact that such analyses leave open how exactly indirect modifiers are structurally integrated in the DP, Japanese does exhibit direct modifiers when applying another criterion: that of \textit{non-predicativity}. Taking the predicate nature as a characteristic of indirect modifiers \citep{bolinger:1967:adj, alex:2007:nounphrase, cinque:2010:adj, panayidou:2013:in}, the lack of predicativity is a characteristic of unambiguously direct, or \textit{direct-only} modifiers. Such modifiers, which have sometimes been noted in the literature (see overview in \citealt{watanabe:2017:attributive}), include \is{non-intersective modifier} non-intersective modifiers such as \textit{shōrai-no} `future’ depicted in \Ref{shorai} and modifiers of \textsc{nationality} and \textsc{material} as depicted in \Ref{denmaku1}.\footnote{In such examples, I depict possible usage in context-neutral sentences, abstracting away from contexts which ameliorate these modifiers in predicative position. These include contrastive contexts \citep{hoshi:2002:kaynean}, contexts in which the head noun is elided in the second instance, thus essentially a case of modification \citep{nagano:2018:conversion}, or contexts in which the copula acts as a placeholder for another verb, the so-called \is{eel-construction}eel-constructions \citep{okutsu:1978:boku}. See Sections \ref{testsdirect} and \ref{non-intersective Japanese}. \label{foonoteeel}} Intriguingly, most direct-only modifiers co-occur with the element \textsc{no}.

\ex. \label{shorai} \ag. shōrai-no daitōryō\\
future-\textsc{no} prime.minister\\
`future prime minister’
\bg. *Kono daitōryō-wa shōrai-da.\\
this prime.minister-\textsc{top} future-\textsc{cop}\\
(`This prime minister is future.’) \\ \null \hfill direct modifier, non-intersective

\ex. \label{denmaku1} \ag. Denmāku-no / chīku-no isu\\
Denmark-\textsc{no} / teak-\textsc{no} chair\\
`a Danish chair / a teak chair'
\bg. *Ano isu-wa Denmāku-da / chīku-da.\\
that chair-\textsc{top} Danish-\textsc{cop} / teak-\textsc{cop}\\
(`This chair is Danish / teak.’) \\ \null \hfill direct modifier, \textsc{nationality/material}

Crucially, even among direct-only modifiers no \is{absence of ordering restrictions}ordering restrictions can be detected in Japanese, contrary to, for instance, English.

\ex. \ag. Denmāku-no chīku-no isu\\
Denmark-\textsc{no} teak-\textsc{no} chair\\
`Danish teak chair' \label{denmaku2}
\bg. chīku-no Denmāku-no isu\\
teak-\textsc{no} Denmark-\textsc{no} chair\\
`teak Danish chair' \label{denmaku3}

\ex. \label{danish english} \a. Danish teak chair
\b. ??teak Danish chair

\newpage
The first main goal of this book is therefore to derive a DP-structure in which both types of modifiers, direct and indirect, can be integrated and which can account for this freedom of ordering in Japanese, which is rather surprising for cartographic accounts \citep{cinque:2010:adj, cinque:2017:micro}. This investigation also includes a precise analysis of how both types of modifiers, direct and indirect, are integrated.

Finally, it is a desideratum to analyze the internal structure of both types of modifiers. First, this concerns morphologically complex elements, for example, \textsc{katta} marking past inflections of \textit{i}-adjectives.

\exg. ano taka-katta doresu\\
that expensive-\textsc{katta} dress\\
`that dress which was expensive' \null \hfill \textit{i}-adjective, past \label{taka-kattaoben}

Second, this concerns the elements consisting only of one mora, namely \textsc{i} following the \textit{i}-adjective \textit{taka-i} `high, expensive’ in \ref{doresuintro}, \textsc{na} following \textit{na}-adjectives, and the most versatile \textsc{no} occurring in \ref{shorai}.\footnote{I will adopt morae as the smallest prosodic units in Japanese, rather than syllables. Cf. for example \citet{kubozono:2005:rendaku, ito:2015:word}.} The internal structure of modifiers is often treated orthogonally to their role inside the DP, but by bringing these issues together, I derive a complete structure of the Japanese DP, also including the micro-level.

\section{Main hypotheses}

This section presents the main hypotheses to be defended in this book. They concern the structure of the DP, the absence of specific functional projections and the internal structure of modifiers.

\subsection{Structure of the DP}

I propose that the Japanese DP looks as depicted in \figref{fig:Structure of the Japanese DP1}.


\begin{figure}

 \begin{subfigure}{0.45\textwidth}

 \small
 \begin{tikzpicture}[baseline=0pt, remember picture]
\tikzset{level distance=30pt, sibling distance=-5pt}
 \Tree [.DP\textsubscript{external} {} [.D′\textsubscript{external} [.FP\textsubscript{QUANTITY} XP [.F´\textsubscript{QUANTITY} [.FocP XP [.Foc′ [.FP\textsubscript{n} XP [.F´\textsubscript{n} [.FP\textsubscript{n+1} XP [.F′\textsubscript{n+1} [.\textbf{DP\textsubscript{internal}} {} ... ] F ] ] F ] ] Foc ] ] F ] ] D ]]
\end{tikzpicture}
 \end{subfigure}
 \begin{subfigure}{0.45\textwidth}

 \begin{tikzpicture}[baseline=0pt, remember picture]
\tikzset{level distance=30pt, sibling distance=-5pt}
 \Tree [.\textbf{DP\textsubscript{internal}} {} [.D′\textsubscript{internal} [.FP\textsubscript{indirect-n} XP [.F´\textsubscript{indirect-n} [.FP\textsubscript{Num} NumP [.F′\textsubscript{Num} [.FP\textsubscript{indirect-n+1} XP [.F′\textsubscript{indirect-n+1} [.\textbf{dP} [.FP\textsubscript{direct-n} XP [.F′\textsubscript{direct-n} [.FP\textsubscript{direct-n+1} XP [.F′\textsubscript{direct-n+1} [.$\sqrt{}$P XP [.$\sqrt{}$′ \edge[roof]; {N} ] ] F ] ] F ] ] d ] F ] ] F ] ] F ] ] D ] ]
\end{tikzpicture}
 \end{subfigure}
 \caption{Structure of the Japanese DP}
 \label{fig:Structure of the Japanese DP1}
\end{figure}


I adopt a split-DP approach \citep{laenzlinger:2005:french, laenzlinger:2011:elements, kim:2019:syntax} and propose that the Japanese DP consists of a DP-external domain, the left periphery, and a DP-internal domain. The DP-internal domain consists furthermore of a domain for indirect (predicative) modification and a lower domain for direct (attributive) modification. Those two domains are separated by the indefinite head d.\footnote{The head d is the small indefinite head of a relative clause if there is one in the structure \citep{cinque:2010:adj, cinque:2020:relcl}. See Section \ref{indefinitedpdiscussion} for discussion.} Note that the DP-external domain is represented by the left part of the tree, and the DP-internal domain is represented by its right part. The thresholds between the three domains are depicted in boldface.

This book focuses on the DP-internal domain, but a discussion on the left periphery can be found in Section \ref{externaldp}, where I show that this area must encompass landing sites for quantificational modifiers, non-restrictive relative clauses and focused modifiers, but that essentially all types of modifiers can move there.

Modifiers, unequivocally depicted as XPs in \figref{fig:Structure of the Japanese DP1}, merge in the specifier position of functional projections (FPs) within the DP. A crucial assumption that differentiates the system defended here from accounts such as \citet{cinque:2010:adj, cinque:2020:relcl} is that functional projections (with one exception) are neither dedicated to specific semantic categories in the direct domain, nor to different sizes of relative clauses in the indirect domain. In other words, these functional projections are unspecific. The labels \textit{n}, \textit{n+1}, are meant to depict that there is potentially an infinite number of functional projections inside every domain, but that there is a hierarchy, depicted by the numbers and the ``plus" symbol.

\subsection{Unspecified functional projections predict free order}

I propose that this is the general DP spine valid in every language, but that cross-linguistic variation should be permitted concerning the nature of the functional projections present in that spine. Some languages, including English, have the relevant language-specific tools, called \is{Units of Language}\textit{Units of Language} in \citet{wiltschko:2014:universal}, available to create functional projections dedicated to fine-grained semantic categories such as \textsc{shape} or \textsc{color}. In turn, this leads to the familiar \is{adjective ordering restrictions}ordering restrictions within the direct domain.

In turn, Japanese modifiers are licensed by functional heads in the respective domains they belong to, but this language lacks the language-specific tools to form \textit{dedicated} functional projections. Since modifiers are not merged in such dedicated, and rigidly ordered, functional projections, their merging position is not predetermined and hence they can occur in any order. As a result, even for a pair of non-predicative modifiers, which should exhibit ordering restrictions given what we know from other languages, no such ordering restrictions are observable.

This is illustrated below in \figref{denmaku vs. chiku} for the two modifiers of \textsc{nationality} \textit{Denmāku-no} `Danish’ and of \textsc{material} \textit{chīku-no} `teak’. Their two possible surface \is{absence of ordering restrictions}orders are simply the result of different merging positions within the DP. See Section \ref{discussiondirect} for discussion on ordering in the direct domain.

% \begin{figure}
%
% \begin{subfigure}{\textwidth}
% \centering
% \begin{tikzpicture}[baseline=0pt, remember picture]
% \tikzset{level distance=25pt, sibling distance=5pt}
% \Tree [.DP\textsubscript{internal} {} [.D′\textsubscript{internal} [.FP\textsubscript{indirect} {} [.F′\textsubscript{indirect} [.dP [.FP\textsubscript{direct} [.XP \edge[roof]; Denmāku-no ] [.F′\textsubscript{direct} [.FP\textsubscript{direct} [.XP \edge[roof]; chīku-no ] [.F′\textsubscript{direct} [.NP \edge[roof]; isu ] F ] ] F ] ] d ] F ] ] D ] ]
% \end{tikzpicture}
% \caption{\textit{Denmāku-no chīku-no isu} (= \ref{denmaku2})}
% \label{fig:denmaku-no-chiku-no}
% \end{subfigure}
%
% \begin{subfigure}{\textwidth}
% \centering
% \begin{tikzpicture}[baseline=0pt, remember picture]
% \tikzset{level distance=25pt, sibling distance=5pt}
% \Tree [.DP\textsubscript{internal} {} [.D′\textsubscript{internal} [.FP\textsubscript{indirect} {} [.F′\textsubscript{indirect} [.dP [.FP\textsubscript{direct} [.XP \edge[roof]; chīku-no ] [.F′\textsubscript{direct} [.FP\textsubscript{direct} [.XP \edge[roof]; Denmāku-no ] [.F′\textsubscript{direct} [.NP \edge[roof]; isu ] F ] ] F ] ] d ] F ] ] D ] ]
% \end{tikzpicture}
% \caption{\textit{chīku-no Denmāku-no isu} (= \ref{denmaku3})}
% \label{fig:chiku-no-denmaku-no}
% \end{subfigure}
%
% \caption{Two possible orders of \textit{Denmāku-no} and \textit{chīku-no}}
% \label{denmaku vs. chiku}
%
% \end{figure}

\begin{figure}
\begin{subfigure}{.45\textwidth}
\centering
\begin{forest}
for tree={
  l sep=25pt,
  s sep=5pt,
}
[DP\textsubscript{internal}
  [{}]
  [D′\textsubscript{internal}
    [FP\textsubscript{indirect}
      [{}]
      [F′\textsubscript{indirect}
        [dP
          [FP\textsubscript{direct}
            [XP
              [Denmāku-no, roof]
            ]
            [F′\textsubscript{direct}
              [FP\textsubscript{direct}
                [XP
                  [chīku-no, roof]
                ]
                [F′\textsubscript{direct}
                  [NP
                    [isu, roof]
                  ]
                  [F]
                ]
              ]
              [F]
            ]
          ]
          [d]
        ]
        [F]
      ]
    ]
    [D]
  ]
]
\end{forest}
\caption{\textit{Denmāku-no chīku-no isu} (= \ref{denmaku2})}
\label{fig:denmaku-no-chiku-no}
\end{subfigure}\hfill
\begin{subfigure}{.45\textwidth}
\centering
\begin{forest}
for tree={
  l sep=25pt,
  s sep=5pt,
}
[DP\textsubscript{internal}
  [{}]
  [D′\textsubscript{internal}
    [FP\textsubscript{indirect}
      [{}]
      [F′\textsubscript{indirect}
        [dP
          [FP\textsubscript{direct}
            [XP
              [chīku-no, roof]
            ]
            [F′\textsubscript{direct}
              [FP\textsubscript{direct}
                [XP
                  [Denmāku-no, roof]
                ]
                [F′\textsubscript{direct}
                  [NP
                    [isu, roof]
                  ]
                  [F]
                ]
              ]
              [F]
            ]
          ]
          [d]
        ]
        [F]
      ]
    ]
    [D]
  ]
]
\end{forest}
\caption{\textit{chīku-no Denmāku-no isu} (= \ref{denmaku3})}
\label{fig:chiku-no-denmaku-no}
\end{subfigure}
\caption{Two possible orders of \textit{Denmāku-no} and \textit{chīku-no}}
\label{denmaku vs. chiku}
\end{figure}

Concerning indirect modifiers, a similar situation holds. They only come in one size in Japanese, namely CP, which means Japanese only has one type of (head-external) relative clause. Again, this predicts the \is{absence of ordering restrictions}absence of ordering restrictions, within the indirect domain as well. However, the indirect domain also encompasses a projection for \is{numeral}numerals and, in principle, indirect modifiers can occur in a position above or below the relevant projection hosting numerals.

\ex. \ag. [pātī-de at-ta [futari-no josei]]\\
party-\textsc{loc} meet-\textsc{pst} two-\textsc{no} woman\\
\bg. [futari-no [pātī-de at-ta josei]]\\
two-\textsc{no} party-\textsc{loc} meet-\textsc{pst} woman\\
`the two women who I met at a party'

Furthermore, indirect modifiers of higher \is{complexity}complexity, for instance when they contain an overt subject or an adverbial, are more natural in a higher position.

\ex. \ag. [doresu-ga utsukushi-k-at-ta [futari-no josei]]\\
dress-\textsc{nom} beautiful-\textsc{pred-cop-pst} two-\textsc{no} woman\\
\bg. \#[futari-no [doresu-ga utsukushi-k-at-ta josei]]\\
two-\textsc{no} dress-\textsc{nom} beautiful-\textsc{pred-cop-pst} woman\\
`the two women whose dresses were beautiful'

This suggests that indirect modifiers originate low, in a position below FP\textsubscript{Num}, but can move up. This discussion will be picked up in Section \Ref{complexity}.

\subsection{Internal structure of modifiers}

The final hypothesis concerns the internal structure of modifiers. Following \citet{nishiyama:1999:adj}, I analyze complex elements such as \textsc{katta}, given in \ref{taka-kattaoben} as combinations of two copulas and a tense marker, as glossed in the following example. This analysis will be outlined in Section \Ref{copulaelements}.

\exg. ano taka-k-at-ta doresu\\
that expensive-\textsc{pred-cop-pst} dress\\
`that dress which was expensive' \null \hfill \textit{i}-adjective, past

On the other hand, the monomoraic elements co-occurring with Japanese non-verbal modifiers in attributive position are not manifestations of copulas and tense markers. Rather, I argue that these elements, exemplified below with \textsc{no}, are overt manifestations of a syntactic head labeled \is{MOD@\textsc{mod}}\textsc{mod} \citep{rubin:1994:mod}. This proposal is similar to the Mod-Insertion hypothesis of \citet{kitagawa:1982:prenominal}, but crucially does not assume post-syntactic insertion. Instead, \textsc{mod} has a syntactic function and licenses the preceding modifier in attributive position along the lines illustrated in \figref{Example internal structure} based on the direct modifier in example \ref{tadanogakusei}.

\exg. kono tada-no gakusei\\
this mere-\textsc{mod} student\\
`this mere student’ \label{tadanogakusei} 

\begin{figure}
\centering
\begin{tikzpicture}[baseline=0pt, remember picture]
\tikzset{level distance=25pt, sibling distance=5pt}
 \Tree [.DP\textsubscript{internal} kono [.D´\textsubscript{internal} [.FP\textsubscript{direct} [.ModP {} [.Mod′ [.XP \edge[roof]; tada ] [.Mod no ] ] ] [.F′\textsubscript{direct} [.NP \edge[roof]; gakusei ] F ] ] D ]]
\end{tikzpicture}
\caption{Example of internal structure of modifiers (= \ref{tadanogakusei})}
\label{Example internal structure}
\end{figure}

I argue for a unified analysis of \is{MOD@\textsc{mod}}\textsc{no} which is not contingent on the nature of the different elements that precede it. Furthermore, the elements following adjectives are instances of such a syntactic head as well, which means its spell-out form is contingent on the word class of the preceding modifier. Readers interested in this discussion are referred to Section \ref{internalstructure}.

\section{The Japanese nominal domain} \label{intro japanese}

In this section, I will present the characteristics of the Japanese language that are relevant for this book. I will focus on the nominal domain and nominal modifiers.

\subsection{General information and word classes}

Japanese is the first language of approximately 125 million people and the official language of Japan and certain regions of Palau. It is a strictly head-final SOV language in which the verb cannot be moved from sentence-final position (\citealt{kuno:1973:structure, martin:1975:reference, iwasaki:2013:japanese} among others). Another key feature is its agglutinative character. Grammatical information is encoded via affixes, predominantly suffixes. \is{word classes}

There is a classical divide between inflecting and uninflecting, as well as dependent and independent word classes,\footnote{Following the terminology of \citet{spencer:2020:uninflectedness}, \textit{inflecting} denotes all those word classes in a specific language whose members generally do inflect. \textit{Uninflecting} denotes word classes whose members never inflect, while \textit{uninflectable} denotes specific members of an otherwise inflecting word class that for various reasons do not inflect.} making \textit{inflection} a key characteristic for differentiating between word classes. What is meant by this term is not the familiar morphological alternation dependent on categories such as number or person; rather, inflection (Japanese: \textit{katsuyō} \textjpn{活用}) denotes the conjugation of verbal categories in different inflectional forms. These inflectional forms, or \textit{katsuyōkei} (\textjpn{活用形}), not only encode different functions, they are also encoded for which functional, or lexical, elements they host.\footnote{Six inflectional forms are recognized on the basis of Classical Japanese: I \textit{mizenkei} \textjpn{未然形} (negative form/irrealis), II \textit{renyōkei} \textjpn{連用形} (nominalization/adverbial form), III \textit{shūshikei} \textjpn{終止形} (sentence-final form, also conclusive), IV \textit{rentaikei} \textjpn{連体形} (attributive form), V \textit{izenkei} \textjpn{已然形} (concessive form) and VI \textit{meireikei} \textjpn{命令形} (imperative form). These forms and their respective terminology are still referred to with regard to Modern Japanese, but the number and nature of these inflectional forms is subject to debate, due to the fact that for many inflectional paradigms, the different forms have become syncretic and hard to distinguish (see \citealt{kohlich:2023:uninflecting} for an overview). Furthermore, as a result of diachronic phonological processes, it is, sometimes, hard to distinguis between inflectional form and suffix, making inflection in Modern Japanese appear almost fusional. Some clear instances of inflectional forms in Modern Japanese can be observed for verbs of the \textit{consonant inflectional paradigm}, for example \textit{nom} `to drink'. Compare the following stem vowel alternations.

\ex. \ag. nom-\textbf{u}\\
drink-\textsc{prs}\\
`I/you/etc. drink' \null \hfill sentence-final
\bg. nom\textbf{i}-ta-i\\
drink-\textsc{vol}-\textsc{prs}\\
`I/you/etc. want to drink' \null \hfill nominalization
\cg. nom\textbf{e}-ba\\
drink-if\\
`if I/you/etc. drink' \null \hfill concessive


Note also that I use the glosses \textsc{shk} (\textit{shūshikei}), \textsc{rtk} (\textit{rentaikei}) and \textsc{ryk} (\textit{renyōkei}) mostly with reference to Classical Japanese and only when relevant with reference to Modern Japanese.} Inflecting word classes include verbs, the two adjective classes, and the functional category of so-called \textit{auxiliary verbs} (Japanese: \textit{jodōshi} \textjpn{助動詞}).\footnote{As will become obvious on multiple occasions, it is a matter of debate whether \textit{na}-adjectives do inflect (see \citealt{uehara:1998:syntactic, kishimoto:2016:lexical}).} \is{Classical Japanese}

\subsection{Modifiers and modification}

The strict head-finality is also apparent in the nominal domain. Modifiers exclusively precede the noun and include numerals, verbs, nominals, two adjective groups, and demonstratives, but not articles. The two adjective groups, in Japanese called \textit{keiyōshi} \textjpn{形容詞}, Lit. `qualifying words’, will consistently be referred to as \textit{\textit{i}-adjectives} and \textit{\textit{na}-adjectives} throughout this book reflecting their systematic morphological differences. \textit{i}-adjectives such as \textit{taka-i} `high, expensive’ end in -\textsc{i} in attributive position, but also in predicative position to mark non-past, whereas \textsc{na}-adjectives such as \textit{shikku-na} `chic’ end in -\textsc{na} in attributive position in non-past. In predicative position, they exhibit the form -\textsc{da}, and this difference essentially reflects the difference between \textit{shūshikei}, the sentence-final form, and \textit{rentaikei}, the attributive form, of Classical Japanese.\footnote{This difference has been reduced in all other inflecting word classes since around Late Middle Japanese (12–16 AD), when the attributive form started to appear in sentence-final environments. The dichotomy is still observable for inflecting word classes in the different varieties of Okinawan (Ryūkyūan), which go back to the same origin as Japanese, but in many instances have preserved features of Proto-Japanese.} \is{Classical Japanese} \is{word classes}

Thus, \textit{na}-adjectives, in non-past, alongside nouns to be reviewed below, provide an exception to the observation that Japanese lexemes always appear in the same form in attributive and predicative position.

\ex. \ag. Ano doresu-wa shikku-da.\\
that dress-\textsc{top} chic-\textsc{da}\\
`That dress is chic.'
\bg. ano shikku-na doresu\\
that chic-\textsc{na} dress\\
`that chic dress' \null \hfill \textit{na}-adjective, non-past

\subsection{Adjectives, I and NA} \label{chapter 1, i and na}

The two adjective groups exhibit the typical differences observable in languages with two adjective groups \citep{dixon:2004:adj}. \textit{i}-adjectives are the older group and most lexemes belong to the native Japanese stratum, whereas \textit{na}-adjectives entered the language in Late Middle Japanese and have only received the status as an independent word class in the middle of the 20th century (\citealt{kato:2003:nihongo, lewin:2003:abriss, frellesvig:2010:history} a.o.). Most lexemes of this adjective class belong to non-native, foreign strata, mostly Sino-Japanese, but also English, but there are also some native lexemes, often with typical endings -\textit{ka}, -\textit{raka} or -\textit{yaka} (see for instance \citealt[210–213]{uehara:1998:syntactic} for diachronic differences). This also means that \textit{na}-adjectives constitute an open class, accepting new loanwords, while \textit{i}-adjectives are a closed class.

As for their semantics, \textit{i}-adjectives typically express basic notions, while more detailed meanings are denoted via \textit{na}-adjectives. However, both groups can in principle express all semantic categories of various well-known hierarchies (\citealt{scott:2002:stacked} among others), although some categories are biased towards one group. More basic categories such as \textsc{dimension} show a higher percentage of \textit{i}-adjectives, while the more versatile category \textsc{quality} incorporates more \textit{na}-adjectives \citep{uehara:1998:syntactic}.

There exist various analyses of the monomoraic elements \textsc{i} and \textsc{na}. Until the 2000s, \textsc{i} was analyzed either as an inflecting tense morpheme \citep{teramura:1982:nihongo, murasugi:1991:noun, urushibara:1993:syntactic, uehara:1998:syntactic}, or (part of) a copula \citep{kubo:1992:japanese}. \textsc{na}, on the other hand, was perceived less controversially and predominantly analyzed as the copula, alongside its \is{predicative copula}predicative equivalent \textsc{da}, which generally received unequivocal copula-status (for example \citealt{urushibara:1993:syntactic, nishiyama:1999:adj}). One reason for this is that not only \textsc{da}, but also the attributive \textsc{na}, can be paraphrased by the complex element \textsc{dearu}, which has been identified as the canonical form of the copula and is historically related to \textsc{da} (see for example \citealt{urushibara:1993:syntactic, nishiyama:1999:adj, nakahara:2002:japanese}; see \citealt{miyama:2011:da} for historical evidence). This is shown below.\footnote{The two elements do not display a significant difference in meaning, but \textsc{dearu} is mostly observed in written register and rarely in spoken language. It is also degraded in attributive position, although not illicit. I will gloss both elements simply as \textsc{cop} for now. For the relevant analysis see Section \ref{copulaelements}.}

\ex. \label{na-A de aru} \ag. Ano doresu-wa shikku-dearu.\\
that dress-\textsc{top} chic-\textsc{cop}\\
`That dress is chic.'
\bg. ano shikku-dearu doresu\\
that chic-\textsc{cop} dress\\
`that dress which/that is chic' \\ \null \hfill \textit{na}-adjective, non-past with copula \textsc{dearu}

\largerpage
Furthermore, both adjective groups can express past tense, and when they do, they exhibit the same complex morphological forms in attributive as in predicative position, meaning that the formerly observed difference in \textit{na}-adjectives disappears. They co-occur with the complex element \textsc{deatta}, which corresponds to \textsc{dearu} in non-past, or the shorter form \textsc{datta}, where the syllable /e/ has been deleted and which corresponds to \textsc{da}.\footnote{The same as \textsc{dearu} and \textsc{da}, these two forms display a difference in register, but are otherwise interchangeable.}

\textit{i}-adjectives on the other hand seem to truly inflect for past. The complex element \textsc{katta} occurs instead of \textsc{i}, again in both attributive and predicative position. At first glance, these two elements seem to bear no resemblance to one another.

\ex. \ag. Ano doresu-wa shikku-d(e)atta.\\
that dress-\textsc{top} chic-\textsc{cop.pst}\\
`That dress was chic.'
\bg. ano shikku-d(e)atta doresu\\
that chic-\textsc{cop.pst} dress\\
`that dress which was chic' \null \hfill \textit{na}-adjective, past

\ex. \label{taka-katta} \ag. Ano doresu-wa taka-katta.\\
that dress-\textsc{top} expensive-\textsc{pst},\\
`That dress was expensive.'
\bg. ano taka-katta doresu\\
that expensive-\textsc{pst} dress\\
`that dress which was expensive' \null \hfill \textit{i}-adjective, past

A landmark analysis to be adopted in this study comes from \citet{nishiyama:1999:adj}, who aimed at a unified analysis of \textsc{i} and \textsc{na} as combinations of two copulas and a tense marker. This analysis, although its adoption at face value brings about some problems, is important in the sense that it offers a parallel analysis of the two elements. In earlier studies, often, but not always \citep{backhouse:1984:have}, the dichotomy between \textit{inflecting tense marker} (\textsc{i}) vs. \textit{copula} (\textsc{na}) reflected the word class status of relevant lexemes, \textit{i}-adjectives were analyzed as (verbal) inflecting predicates, and \textit{na}-adjectives as non-inflecting nouns (see most of all \citealt{uehara:1998:syntactic}).

What should be pointed out, however, is that the modifiers that co-occur with \textsc{i} and \textsc{na} are adjectives.\footnote{The fact that \textit{i}-adjectives seem rather ``verby" and \textit{na}-adjectives rather ``nouny" has prompted several researchers to subsume them under the headlines of \textit{Verbals} and \textit{Nominals} \citep{uehara:1998:syntactic}, sometimes denying the adjective category in Japanese in general \citep{wetzer:1996:typology}. In fact, the analysis of \citet{nishiyama:1999:adj} and its subsequent incorporation in the study of \citet{baker:2003:cat} were important steps towards the assumption of two adjective groups in Japanese.} First, note that, abstracting away from semantic factors, such lexemes are typically gradable as exemplified for a \textit{na}-adjective in the following example. Gradability is a typical feature of the adjectival category \citep{baker:2003:cat, dixon:2010:grammatical}, that in Japanese is expressed via comparative-formation with the particle \textit{yori}(-\textit{mo)} `than' and the optional adverb \textit{motto} `more' \ref{machi-yori}, degree adverbs \ref{totemonigiyaka}, or superlative-formation via the adverb \textit{ichi-ban} `most' (Lit. `one-number') \ref{machi-ichiban}, see \citet{bhatt:2011:reduced}.

\ex. \ag. Kono machi-wa ano machi-yori (motto) nigiyaka-da.\\
this city-\textsc{top} that city-than (more) lively-\textsc{cop}\\
`This city is more lively than that city.’ \label{machi-yori}
\bg. totemo nigiyaka-na machi\\
very lively-\textsc{na} city\\
`very lively city' \label{totemonigiyaka}
\cg. ichi-ban nigiyaka-na machi\\
number-one lively-\textsc{na} city\\
`the most lively city' \label{machi-ichiban}

Another, more language-specific, feature of adjectivehood is \is{nominalization}nominalization with the suffix -\textit{sa} `-ness’, which only nominalizes adjectives, but never nouns, since they are already nominal. As shown below, \textit{i}-adjectives (\textit{samu-i} `cold’) and \textit{na}-adjectives (\textit{shizuka-na} `quiet’) can be nominalized, but nouns (\textit{sensei} `teacher’) or verbs (\textit{tabe} `eat’) cannot.\footnote{See \citet{teramura:1982:nihongo, wenck:1974:systematische, urushibara:1993:syntactic, rickmeyer:1995:japanische, kato:2003:nihongo, muraki:2012:nihongo} for a similar application of this characteristic. Other, more specific suffixes which only attach to adjectives are -\textit{garu}, which derives verbs from adjectives that express (often negative) physiological or psychological states \citep{urushibara:1993:syntactic}, and the verbal suffix -\textit{sugiru} `to exceed; too’.}

\ex. \ag. samu-sa\\
cold-ness\\
\bg. shizuka-sa\\
quiet-ness\\
\cg. *sensei-sa\\
teacher-ness\\
\dg. *taberu-sa\\
eat-ness\\

Another characteristic differentiating adjectives from nouns is the inability to appear with particles, especially -\textit{ga} for nominative and -\textit{wo} for accusative case \citep{suzuki:1972:nihon, martin:1975:reference, teramura:1982:nihongo, nakayama:2004:tokubetsu, muraki:2012:nihongo}, in other words the lack of \textit{referential use} \citep{oshima:2019:gradability}, as shown below.

\ex. \ag. *Samu-(i-)ga ku-ru.\\
cold-\textsc{i-nom} come-\textsc{prs}\\
(`(The/An) cold comes’.)
\bg. *Samu-(i-)wo mi-ru.\\
unknown-\textsc{i-acc} see-\textsc{prs}\\
(`to see (the/a) cold’) \null \hfill \textit{i}-adjective

\ex. \ag. *Shizuka-ga ku-ru.\\
quiet-\textsc{nom} come-\textsc{prs}\\
(`(The/An) quiet comes.’)
\bg. *Shizuka-wo mi-ru.\\
quiet-\textsc{acc} see-\textsc{prs}\\
(`to see (the/a) quiet’) \null \hfill \textit{na}-adjective

\subsection{\textit{no}-Modifiers and \textsc{NO}}

The discussion above was meant to show that \textsc{i} and \textsc{na} co-occur only with adjectives. This subsection discusses the monomoraic element \textsc{no}, whose distribution is not restricted to a specific word class.

\subsubsection{Nouns}

In most grammars of Japanese, \textsc{no} is strongly, often exclusively, connected to its co-occurrence with nouns. In fact, nouns share their entire inflectional paradigms with \textit{na}-adjectives, except for this element (\textsc{no}) which they co-occur with in attributive position, contrary to the element \textsc{na} that co-occurs with \textit{na}-adjectives. The prototypical relationship between two nouns connected via \textsc{no} is arguably possession, which explains why this element is predominantly glossed as \textsc{gen}.\footnote{For discussion and different opinions see \citet{murasugi:1991:noun, uehara:1998:syntactic, backhouse:2004:inflected}. I will gloss it as \textsc{no} until I present the proper analysis as a syntactic head in Section \ref{internalstructure}.}

\ex. \label{1:senseinohon} \ag. sensei-no hon\\
teacher-\textsc{no} book\\
`book of the teacher'
\bg. Tarō-no jitensha\\
Taro-\textsc{no} bicycle\\
`Taro’s bicycle’ \null \hfill noun, possessive

The lexemes combined via \textsc{no} in \ref{1:senseinohon} are clearly nouns. Syntactically, they can appear with case particles as in \ref{noun reference}; from a semantic and pragmatic point of view, they establish a \textit{reference} to some entity, most obvious with the proper name \textit{Tarō}, alongside common ground between speaker and hearer.

\exg. sensei-wo / Tarō-wo / jitensha-wo / hon-wo mi-ta\\
teacher-\textsc{acc} / Taro-\textsc{acc} / bicycle-\textsc{acc} / book-\textsc{acc} see-\textsc{pst}\\
`saw the teacher / Taro / the/a bicycle / the/a book’ \label{noun reference}

Several other relationships can be invoked when two nouns are connected via \textsc{no}, prominently argument roles \citep{murasugi:1991:noun}. The following example is ambiguous between a subject and an object reading.

\exg. Itaria-no hakai\\
Italy-\textsc{no} destruction\\
1. `Italian destruction' \null \hfill (= destruction by Italy, \textit{subject})\\
2. `Italian destruction' \null \hfill (= destruction of Italy, \textit{object})

A subject and an object can even appear side by side within the same DP, both of them followed by \textsc{no}. Usually (abstracting away from meaning), the reading where the first noun is interpreted as the subject is more accessible, but permutation is naturally possible.\footnote{Examples of multiple occurrences of \textsc{no} also show that it is recursive in the sense that in theory an infinite number of elements can be connected via \textsc{no} in the same DP. This is similar to English \textit{of} or genitive -'\textit{s}, with the difference that two or more argument genitives are not allowed in one DP as also visible in the translations of the relevant example \citep[86]{murasugi:1991:noun}.}

\exg. yabanjin-no toshi-no hakai\\
barbarian-\textsc{no} city-\textsc{no} destruction\\
`the destruction of the city by the barbarians’ \hfill \citep[7]{murasugi:1991:noun}

These DPs can be paraphrased as sentences where the modifying nouns bear nominative and accusative case particles and the modified noun forms a predicate with the light verb \textit{su} (inflected for \textit{renyōkei} in \ref{yabanjin}), also showing that \textsc{no} is mutually exclusive with case particles.

\exg. Yabanjin-ga toshi-wo hakai-shi-ta.\\
barbarian-\textsc{nom} city-\textsc{acc} destruction-do-\textsc{pst}\\
`The barbarians destroyed the city.’ \label{yabanjin}

\largerpage
Furthermore, \textsc{no} can license the attributive modification of nouns that bear other diverse argument roles. In this case, those nouns typically appear with postpositions in order to make the role they bear clearer. These postpositions include -\textit{de} (instrumental, locative), -\textit{ni} (dative, locative, allative, illative) and -\textit{kara} (ablative) \citep{murasugi:1991:noun,saito:2008:N}.\footnote{When the context is sufficiently clear, the postpositions preceding \textsc{no} can also be omitted.} In addition, note that these examples give a first indication that not everything that co-occurs with \textsc{no} is a noun.\footnote{In fact, this is one argument for the distinction which is often drawn between case particles and postpositions. Case particles do not co-occur with \textsc{no}. Note furthermore that PPs can also co-occur with case particles, which modifiers belonging to the adjective groups can never do \citep{murasugi:1991:noun, urushibara:1993:syntactic}.}

 \ex. \label{no with case} \ag. ishi-de-no kōgeki\\
stone-\textsc{ins-no} attack\\
`an attack with stones' \null \hfill instrumental \citep[248]{saito:2008:N}
\bg. Tōkyō-de-no kaigi\\
Tokyo-\textsc{loc-no} meeting\\
`the meeting in Tokyo \null \hfill locative
\cg. Tōkyō-kara-no densha\\
Tokyo-\textsc{abl-no} train\\
`the train from Tokyo' \null \hfill ablative \citep[50]{murasugi:1991:noun}

Note that \textsc{no} always follows the respective postposition and must be closest to the noun, highlighting its connecting function.

\exg. *ishi-no-de kōgeki\\
stone-\textsc{no-ins} attack\\

Nouns acting as modifiers do not necessarily have to bear an argument status. Sometimes they are adjuncts \ref{adjunct no}. This also includes cases of possessor-possessee relationship where the possessor appears as the second noun \ref{alienable}.

 \ex. \label{adjunct no} \ag. ame-no hi\\
rain-\textsc{no} day\\
`rainy day' \null \hfill \citep[251]{saito:2008:N}
\cg. otoko-no ko\\
man-\textsc{no} child\\
`male child (=boy)'
\cg. kuruma-no zasshi\\
car-\textsc{no} magazine\\
`magazine about cars' \null \hfill adjuncts

\ex. \label{alienable} \ag. megane-no gakusei\\
glasses-\textsc{no} student\\
`the student with glasses'
\bg. shiraga-no gakusei\\
white.hair-\textsc{no} student\\
`the student with white hair' \\ \null \hfill \citep[96]{kato:2003:nihongo}

\subsubsection{Adjectival modification}

Although most Japanese grammars (\citealt{iwasaki:2013:japanese, tsujimura:2014:intro, hasegawa:2015:japanese} etc.), including the so-called \textit{school grammar} (\citealt{hashimoto:1948:kokugogaku}), treat lexemes preceding \textsc{no} as nouns, \textsc{no} is not restricted to the lexical class of nominals, but is also involved in (qualitative) adjectival modification \citep{muraki:2012:nihongo, oshima:2019:gradability}. A good example is offered with the lexeme \textit{mumei-no} `unknown’, which appears (almost) exclusively with \textsc{no}, but exhibits clear characteristics of adjectivehood. These include gradability, at least with adverbs such as \textit{kanzenni} `completely’ \ref{kanzennimumei} and inside comparative clauses \ref{yorimumei}, the possibility to be nominalized via the suffix -\textit{sa} `-ness’ \ref{mumeisa}, and the impossibility to appear with nominative and accusative case particles, \ref{mumei-ga} and \ref{mumei-wo}. For these reasons, such lexemes are sometimes referred to as \is{no-adjective@\textit{no}-adjective}\textit{no-adjectives} \citep{mio:1942:hanashi, mio:1958:hanashi, muraki:2012:nihongo}, making them a third morphological subgroup of the adjectival category in Japanese (for different analyses of \textit{mumei-no} see \citealt{teramura:1982:nihongo, kato:2003:nihongo, morita:2013:morph}; for an overview see \citealt{kohlich:2023:suppl}).

\ex. \ag. mumei-no haiyū\\
unknown-\textsc{no} actor\\
`unknown actor'
\bg. kanzenni mumei-no haiyū\\
completely unknown-\textsc{no} actor\\
`completely unknown actor' \label{kanzennimumei}
\cg. Tanaka-san-yori (motto) mumei-no haiyū-to at-ta.\\
Tanaka-Mr.-than more unknown-\textsc{no} actor-\textsc{com} meet-\textsc{pst}\\
`I met with a more unknown actor than Mr. Tanaka.’ \label{yorimumei}
\dg. mumei-sa\\
unknown-ness\\
`unknownness; anonymity’ \label{mumeisa}
\eg. *Mumei-ga ku-ru.\\
unknown-\textsc{nom} come-\textsc{prs}\\
(`(The/An) Unknown comes.’) \label{mumei-ga}
\fg. *Mumei-wo mi-ru.\\
unknown-\textsc{acc} see-\textsc{prs}\\
(`I see (the/an) unknown.’) \label{mumei-wo}

Importantly, lexical modifiers co-occurring with \textsc{no} are not only clear nouns or adjectives, but include many in-between cases. Some modifiers can occur with \textsc{na} in addition to \textsc{no}, some can occur with both and exhibit nominal use, in addition to some which only occur with \textsc{na}, but are simultaneously nominals.\footnote{For more on this see Section \ref{corpusanalysis} below and \citet{kohlich:2023:suppl}.} I will refer to all modifiers that appear with \textsc{no} as \textit{\textit{no}-modifiers} adopting a uniform non-committal label for relevant elements.

Such lexical \is{no-adjective@\textit{no}-adjective}\textit{no}-modifiers truly share the rest of their inflectional paradigm with \textit{na}-adjectives, meaning they also appear predicatively with the copula \textsc{da}. Both, \textsc{no} and \textsc{da}, can be paraphrased by the canonical copula \textsc{dearu}, as was shown for the parallel examples with \textit{na}-adjectives in \ref{na-A de aru}.

\ex. \ag. Ano haiyū-wa mumei-da.\\
that actor-\textsc{top} unknown-\textsc{cop}\\
`That actor is unknown.’
\bg. Ano haiyū-wa mumei-dearu.\\
that actor-\textsc{top} unknown-\textsc{cop}\\
`That actor is unknown.’
\cg. ano mumei-dearu haiyū\\
that unknown-\textsc{cop} actor\\
`that unknown actor’

\subsubsection{Functional elements} \label{grammatical no}

Besides lexical elements, all functional elements which serve as nominal modifiers are connected to the head noun via \textsc{no}. This holds for most \is{quantificational modifier}quantifiers, which are possibly an in-between case between lexical and functional modifiers.\footnote{Eventually, I will treat quantifiers as lexical elements. See Section \ref{quantificational modifiers}.}

\exg. takusan-no / subete-no gakusei\\
a.lot-\textsc{no} / all-\textsc{no} student\\
`many / all students' \hfill \null quantifiers

Clear functional elements that can appear as attributive modifiers in Japanese are \textit{wh}-elements \ref{wh no} (see for instance \citealt{watanabe:2001:wh}) and classifiers, which follow numerals \ref{classifier no} (\citealt{saito:2008:N}).

\exg. Doko-no ATM-de shukkin deki-mas-u ka.\\
where-\textsc{no} ATM-\textsc{loc} withdraw can-\textsc{aux.pol-prs} \textsc{q}\\
`At which ATMs can you withdraw?’ (Lit. `At ATMs where...’) \label{wh no} \footnote{\url{https://help.jibunbank.co.jp/faq_detail.html?id=182, 2023/02/05}.} \\ \null \hfill \textit{wh}-elements

 \ex. \label{classifier no} \ag. san-bon-no bīru\\
three-\textsc{cl-no} beer\\
`three bottles of beer’ \label{sanbon}
\bg. go-nin-no gakusei\\
five-\textsc{cl-no} student\\
`five students’ \null \hfill numeral classifiers

\textsc{no} is also the final mora of demonstratives of the type \textit{kono} `this’, but here has a bound character in Modern Japanese.

\exg. kono / sono / ano / dono hon\\
this.\textsc{speaker-proximate} / that.\textsc{hearer-proximate} / that.\textsc{distal} / which book\\
`this (proximal) / that (medial)  / that (distal) / which book'

\subsubsection{Analysis of NO}

Compared to \textsc{i} and \textsc{na}, the status of \textsc{no} is more versatile and harder to pinpoint. The question is whether every instance of \textsc{no} can be reduced to a single function, at least DP-internally, or whether different instances of \textsc{no} should be postulated. Possible analyses of \textsc{no} are outlined in detail in Section \ref{comparisonpreviousno}, but this book attempts at a unified analysis of \textsc{no}, namely as a syntactic head that licenses modification. The analysis depicted in \figref{Example internal structure} can therefore be applied to all examples mentioned in this section.\footnote{One exception is the demonstratives of the type \textit{kono} `this’ which will be treated as specifiers of D-heads reflecting standard assumptions in syntactic research on Japanese.} For instance a \is{numeral}numeral (containing a classifier) would be integrated as in \figref{Structure Sanbon-no}.

\begin{figure}
\centering
\begin{tikzpicture}[baseline=0pt, remember picture]
\tikzset{level distance=25pt, sibling distance=5pt}
 \Tree [.DP\textsubscript{internal} {} [.D′\textsubscript{internal} [.FP\textsubscript{Num} [.ModP {} [.Mod′ [.ClP {} [.Cl′ [.NumP \edge[roof]; {san} ] [.Cl bon ] ] ] [.Mod no ] ] ] [.F´\textsubscript{Num} [.NP \edge[roof]; bīru ] F ] ] D ]]
\end{tikzpicture}
\caption{\textit{san-bon-no bīru} `three bottles of beer’ (= \ref{sanbon})}
\label{Structure Sanbon-no}
\end{figure}

\subsubsection{Where \textsc{NO} is not used} \label{no no}

It seems that \textsc{no} appears in all instances of (non-adjectival) modification, which is why it is often treated similarly to Mandarin \textsc{de} (see \citealt{saito:2008:N}). However, it is illicit after verbal relative clauses including a finite (past or non-past) verb (see for instance \citealt{murasugi:1991:noun, saito:2008:N}).

\ex. \ag. kare-ga kat-ta (*no) hon\\
he-\textsc{nom} buy-\textsc{pst} \textsc{no} book\\
`the book that he bought’
\bg. hon-wo yon-de-i-ru (*no) gakusei\\
book-\textsc{acc} read-\textsc{ger}-be-\textsc{prs} \textsc{no} student\\
`the student who is reading a book’

Similarly, \textsc{no} never appears after complex morphological past forms, which also contain a tensed form.

\exg. ano shizuka-datta (*no) yoru\\
that quiet-\textsc{cop.pst} \textsc{no} night\\
`that night which was quiet’

Furthermore, it is never the case that any of the elements \textsc{no}, \textsc{na} and \textsc{i} co-occur.

\exg. aka-i (*no) hana\\
red-\textsc{i} \textsc{no} flower\\
`red flower'

\exg. kanzen-na (*no) baka\\
complete-\textsc{na} \textsc{no} idiot\\
`complete idiot'

\largerpage[2]
Finally, \textsc{no} is never possible in predicative position.\footnote{The form \textit{no-da} exists, however in this case \textit{no} is an emphatic element whose usage domain is pragmatically restricted.}

\exg. Ano haiyū-wa mumei (*no) da.\\
that actor-\textsc{top} unknown \textsc{no} \textsc{cop}\\
`That actor is unknown.’

\clearpage
\section{Data collection} \label{corpusanalysis}

The theoretical insights of this book are mainly based on an empirical analysis in the \is{BCCWJ}Balanced Corpus of Contemporary Written Japanese (BCCWJ).\footnote{\url{https://chunagon.ninjal.ac.jp/auth/login} \citep{maekawa:2014:bccwj}.} All relevant materials are saved at the following repository: \url{https://osf.io/u9txp/}. I have analyzed approximately 2,300 lexemes which were predominantly collected from the scientific literature, in total from a body of 23 different studies. The categories according to which I incorporated these lexemes were either \textit{(alleged) inability to be used predicatively} or \textit{use with \textsc{no}, but no clear nominal status}. This list of lexemes will be referred to as the \textit{data set}. The repository contains two files with different annotations of this data set. Data set 1 contains a list of all relevant lexemes annotated based on the judgment of each respective author for cross-referencing purposes. Data set 2 displays the frequency of each relevant variable tested for all lexemes in the corpus.\footnote{Compilation of these lexemes ended at the beginning of 2022. Subsequent analysis began in March 2022 and was ended in August 2023. Readers are referred to a research report published by the Tokyo Gakugei University at the beginning of 2022 \citep{abe:2022:no} in which all lexemes given in \citet{muraki:2012:nihongo} are analyzed for the variables \textit{attributive use, predicative use, adverbial use} and \textit{referentiality}. This is therefore a major contribution to the field and many aspects are shared between this research report and my own study although sometimes numbers diverge slightly. See the data paper for relevant details.}

The repository also contains a data paper, which in detail outlines the design and structure of the corpus and its analysis tools and depicts all necessary steps and search queries that were performed \citep{kohlich:2023:suppl}. It also gives a more detailed overview of the results of the analysis that can be granted here. Readers are welcome to replicate this study, verify the results, implement refinements and provide feedback.

My goal was to establish an overview of how big the inventory of direct, that is, non-predicative, modifiers in Japanese can actually be. In this respect, one must be aware that the existing literature on this matter can only be a reference point. Approximately 75\% of the lexemes in the data set are taken from the paper collection published as a monograph by MURAKI Shinjirō \citep{muraki:2012:nihongo}. Muraki presents long (yet non-exhaustive) lists of different kinds of \textit{no}-modifiers drawn from various Japanese dictionaries, among those an important list of non-predicative lexemes \citep[96–98]{muraki:2012:nihongo}, and analyzes them according to his own introspective judgment.\footnote{The validity of introspective judgments, even with regard to seemingly well-known examples and languages, has been put into question in recent years \citep{gibson:2010:weak}, and this also applies to Japanese.} The present corpus analysis had the goal of verifying such judgments.

I therefore tested for the variable \textit{predicative use with copula} for all modifiers in the data set. A modifier was classified as direct (group 5) if it was either completely lacking predicative use, or if it yielded no more than 5 predicative occurrences vs. at least 20 attributive ones, or alternatively a ratio of 5:1 towards attributive use.\footnote{As a rule of thumb, only lexemes with at least 5 occurrences are reported in the data set, and only lexemes with at least 25 occurrences were categorized into one of the different groups. I am aware of the fact that fewer predicative occurrences do not automatically mean that a modifier is non-predicative, in other words direct-only. Nevertheless, a skew towards attributive use signals that this is the prototypical usage domain of a given modifier, which is why I classified relevant lexemes into this group even though they might be possible in predicative position.} The relevant numbers are given below.

\ex. \a. Direct modifiers in total: 511
\b. Out of which direct modifiers with 0 hits in predicative position: 110
\c. Out of which direct modifiers with no more than 5 hits in predicative position: 190
\d. Prototypical direct modifiers with a ratio of 5:1 towards attributive use: 211

Besides the goal of determining which lexemes are candidates for direct modifiers, another important goal was to point out the heterogeneous word class status of \textit{no}-modifiers. To this purpose, the following variables were investigated for each lexeme:\footnote{Consult Appendix B for more information and numbers on all relevant groups.}

\ex. \a. occurrences in attributive use
\b. element in attributive position: \textsc{na} or \textsc{no}
\c. referential use with case particles for subject and object
\d. nominalization

One might object that corpus data do not tell the whole story for a given phenomenon and indeed, corpora depict a representative sample of real-life language and only show what is \textit{attested}, not what is \textit{possible} or \textit{impossible}. When basing our analyses on corpus queries alone we must keep in mind that certain occurrences might be missed or even overlooked. Nevertheless, a corpus is a powerful tool to give us an overview of a language and its patterns, which in turn allows us to draw important generalizations -- if we know how to look for them. As noted in \citet[7]{desagulier:2017:corpus}, ``[E]ven if no corpus can provide access to the true, unknown law of a language, a corpus can be considered a sample drawn from this law".

In order to further balance out these aspects, I have also performed qualitative interviews with native speakers on several occasions and had discussions with many consultants. The discussion on direct modifiers in Chapter \ref{direct japanese} is based on a qualitative study with five native speakers to inquire about their judgments on grammaticality and acceptability. The study on the position of relative clauses in Chapter \ref{indirect japanese} is also based on interviews with three different native speakers. All these measures were undertaken to ensure the empirical adequacy of the data and the correct theoretical implications to be drawn from them.

